\section{Discussion}
\label{sec:discussion}

\paragraph{What the results show.}
CLDB-RL is designed so that retention and adaptation are not in conflict. The
benchmark deliberately makes phases conflict, so a policy that simply fine-tunes
on the most recent phase should collapse; we observe exactly that for the
replay-free control, whose performance on earlier phases decays to the KEEP
baseline (forgetting \(\approx 0.22\)). Restoring a replay buffer recovers
retention almost completely, and the database-aware prioritization and EWC
provide a further, smaller improvement. This ordering---online $\ll$ replay
$\approx$ full---is the central empirical claim.

\paragraph{Why persistence, not just replay, matters.}
Experience replay works here because the conflicting phases share observations:
the buffer keeps a set of examples from every phase, so gradient steps on the
current phase are interleaved with steps that protect the past. This is the
mechanism identified by \citet{rolnick2019replay}. EWC contributes a
complementary, parameter-space force that is useful precisely for the phases
whose examples are no longer in the buffer.

\paragraph{The adapt/retain trade-off in practice.}
Proposition~\ref{prop:drift} predicts that a larger \(\lambda\) reduces drift
(and forgetting) but slows adaptation. In our runs, EWC's benefit is real but
modest, and at short training budgets it can slightly slow re-learning of a
returning phase. This is the expected shape of the trade-off and suggests that
\(\lambda\) should be tuned per deployment---or made adaptive, which is a natural
extension.

\paragraph{Cost.}
Drift-triggered consolidation incurs a Fisher computation only when the workload
changes; the online pipeline runs all learning off the critical path, so query
latency is unaffected by training. The dominant runtime cost of the full model
over the ablations is the per-step evaluation of the EWC penalty, which grows
with the number of consolidated phases; this is a known cost of regularization
methods and can be bounded by consolidating only the most recent $k$ phases or by
using a Kronecker-factored Fisher.

\subsection{Threats to validity}
\label{sec:threats}
\paragraph{External validity.}
The most important limitation is that the benchmark is a \emph{cost model}, not a
running DBMS. The latency model \eqref{eq:latency} abstracts away plan-cost
non-linearities, concurrency, buffer-pool effects, and the interaction of
multiple indexes. A real evaluation on a system such as PostgreSQL with a
standard workload (TPC-H, TPC-C, or the Join Order Benchmark) and measured
latency as the reward is the essential next step; the pipeline already measures
latency when a database connection is available, so the change is confined to the
environment.

\paragraph{Construct validity.}
We measure \emph{reward} and the CL metrics derived from it, not downstream
metrics such as end-to-end throughput or p99 latency. Reward is defined to be
proportional to the operating cost, so this is a faithful proxy, but it is a
proxy.

\paragraph{Statistical validity.}
We use eight seeds and paired tests. The replay-vs-online effect is large and
consistent; the EWC-vs-no-EWC effect is smaller and, at short training budgets,
not always significant. We therefore report EWC's benefit as suggestive rather
than established, and we are explicit that the headline result is the robustness
conferred by replay.

\paragraph{Comparability.}
The baselines are internal (KEEP/RANDOM/GREEDY-CREATE). A rigorous study must
compare against published index advisors and learned optimizers
\citep{chaudhuri1997autoadmin,valentin2000db2advisor,marcus2021bao} and against
RL index selectors, on a shared benchmark. We view CLDB as complementary to those
systems: any one of them could be substituted as the base policy.

\paragraph{Scope of the RL claim.}
The agent is value-based and off-policy (DQN). We do not claim a new RL
algorithm; the contribution is the continual formulation and its
database-specific instantiation. Likewise, the safety mask simplifies
exploration but means the agent is never asked to learn that structurally invalid
actions are bad---by design, since such exploration is unsafe.
